% Generated from common/cover-notes.json by tools/series_cover.py.
\backmatter
\cleardoublepage
\chapter*{About the Cover}
\addcontentsline{toc}{chapter}{About the Cover}
\noindent A curve that bends more sharply the farther it has travelled, with the circle that fits it best at each of several points. Leonhard Euler found the curve in 1744, as the shape of a thin elastic spring. The integrals that define it are the Fresnel integrals, and Alfred Cornu drew the curve in 1874 to study diffraction, which is why it is also called the Cornu spiral. It is a calculation one cannot finish in elementary functions. It belongs with the chapter Integration Techniques, and with the Fourier transform later in the volume. The drawing was computed by the book's own TikZ cover generator.
